\section{\label{sec:match}匹配核的类 DGLAP 方程}

忽略夸克味之间及夸克与胶子之间的混合（即考察非单态分布），重正化 PDFs 满足描述其
能标依赖性的 DGLAP 方程
\cite{Gribov:1972rt,Dokshitzer:1977sg,Altarelli:1977zs}
\begin{equation}
\mu\frac{\mathrm{d}\, f(x,\mu)}{\mathrm{d}\mu} = \frac{\alpha_s(\mu)}{\pi}
\bigintsss_x^1
\frac{\mathrm{d}y}{y} f(y,\mu) P\left(\frac{x}{y}\right).
\end{equation}
这里 $P\left(z\right)$ 是一个函数，其矩由
\begin{equation}\label{eq:gamman}
\int_0^1 \mathrm{d}x\, x^{n-1} P(x) = \gamma^{(n)},
\end{equation}
给出，其中
\begin{equation}\label{eq:rgan}
\left[\mu\frac{\mathrm{d}\;}{\mathrm{d}\mu} -\frac{\alpha_s(\mu)}{\pi}
\gamma^{(n)} \right]a^{(n)}(\mu) = 0,
\end{equation}
而 $\alpha_s(\mu)$ 是（重正化的）强耦合常数。

类似地，我们可以为联系涂抹准 PDFs 与光前 PDFs 的匹配核推导出一个类 DGLAP 方程。从式~\eqref{eq:double} 中的短距离展开出发，应用
重正化群方程 \eqref{eq:rgan}，可得短距离系数满足
\begin{equation}
\left[\mu\frac{\mathrm{d}\;}{\mathrm{d}\mu} +
\frac{\alpha_s(\mu)}{\pi}\gamma^{(n)}\right]C_n^{(0)}(\sqrt{\tau}\mu,\sqrt{\tau}
P_z ) = 0+
{\cal
O}(\sqrt{\tau}\Lambda_{\mathrm{QCD}}),
\end{equation}
及其逆满足
\begin{equation}
\left[\mu\frac{\mathrm{d}\;}{\mathrm{d}\mu}-
\frac{\alpha_s(\mu)}{\pi}\gamma^{(n)}\right]
\left[C_n^{(0)}(\sqrt{\tau}\mu,\sqrt{\tau}
P_z)\right]^{-1} = 0+ {\cal
O}(\sqrt{\tau}\Lambda_{\mathrm{QCD}}).
\end{equation}

把式~\eqref{eq:cinvz} 与 \eqref{eq:gamman} 代入这一重正化群方程，便可得到匹配核的类 DGLAP 方程：
\begin{equation}\label{eq:dglapz1}
 \mu\frac{d\;}{d\mu}\,Z\left(x,\sqrt{\tau}\mu, \sqrt{\tau}P_z\right) =
\frac{\alpha_s(\mu)}{\pi} \bigintsss_{x}^\infty \frac{dy}{y}
Z\left(y,\sqrt{\tau}\mu, \sqrt{\tau}P_z\right) P\left(\frac{x}{y}\right),
\end{equation}
以及
\begin{equation}
 \mu\frac{d\;}{d\mu}\,\widetilde{Z}\left(x,\sqrt{\tau}\mu, \sqrt{\tau}P_z\right) =  -
\frac{\alpha_s(\mu)}{\pi} \bigintsss_{x}^\infty \frac{dy}{y}
\widetilde{Z}\left(y,\sqrt{\tau}\mu, \sqrt{\tau}P_z\right) P\left(\frac{x}{y}\right),
\end{equation}
其误差为 ${\cal O}(\sqrt{\tau}\Lambda_{\mathrm{QCD}})$ 量级。


